\documentclass[10pt,a4paper]{amsart}
\usepackage[margin=19mm]{geometry}
\usepackage{amsmath,amssymb,mathtools}
\usepackage[scr=rsfs,cal=euler]{mathalfa}
\usepackage{xfrac}
\usepackage[unicode,hidelinks,bookmarks=false]{hyperref}
\newcommand{\ie}{i.e.\ }
\newcommand{\cf}{cf.\ }
\newcommand{\resp}{respectively\ }
\newcommand{\Subsection}{\subsection}
\providecommand{\nobreakdash}{\nobreak\hbox{-}}
\setlength{\emergencystretch}{2em}
\multlinegap0pt
\usepackage{bm}
\usepackage{bbm}
\usepackage{dsfont}
\usepackage{xspace}
\usepackage{cancel}
\usepackage{mathtools}
\usepackage{amsthm}
\makeatletter
\@ifundefined{definition}{\newtheorem{definition}{Definition}}{}
\@ifundefined{prop}{\newtheorem{prop}[definition]{Proposition}}{}
\makeatother
\def\tl{\widetilde}
\def\wh{\widehat}
\newcommand{\Z}{\mathbb Z}
\title[Isomoprhism of generalized Bratteli diagrams]{Isomoprhism of generalized Bratteli diagrams}
\author{Olena Karpel}
\date{Source: arXiv:2605.18278v1 (CC BY 4.0)}
\begin{document}
\maketitle

\noindent\emph{Source and licence.} Isomoprhism of generalized Bratteli diagrams. Version arXiv:2605.18278v1 (CC BY 4.0), distributed under the Creative Commons Attribution 4.0 International licence, as recorded in the arXiv licence metadata for this exact version and shown on the abstract page https://arxiv.org/abs/2605.18278v1. Full rights notice: licenses/arxiv-2605.18278-CC-BY-4.0.txt.\par\smallskip

\noindent\textit{Editorial scope.} Counted unit: the Prop whose source label is \texttt{None} in the article (source0.tex lines 1324--1326, source label \texttt{None}), delivered together with the article's own complete original proof of it (source0.tex lines 1328--1358, one \texttt{proof} environment of 31 lines). No mathematics is written, repaired or re-derived: statement and proof are the article's own text line for line. Required context retained uncounted: none. Every definition, display and label of those lines is retained unchanged. Omitted from this package and not redistributed: 1--1323, 1327--1327, 1359--2208. Nothing inside a retained excerpt is omitted or altered: every retained body is delivered exactly as the article's lines are written, so no omission receipt of any kind accompanies this package. Citations: the retained excerpts carry no citation command at all, so no bibliography entry of the article is transcribed and no reference is redistributed. Reference adaptation: none was needed and none was made; every reference inside the delivered unit resolves inside it, so no unresolved reference or citation is produced. Printed slips kept literal: no printed word, symbol, hypothesis or number of the counted statement or of its proof was altered. Layout and class: the article's own class is \texttt{amsart} with packages including inputenc, color, amsthm, amssymb, latexsym, amsmath, fontenc, graphs, MnSymbol, bbm, centernot, mathrsfs, amsfonts, textcomp; the wrapper is the lane's pinned \texttt{amsart} editorial wrapper at 10pt on A4, which restates the theorem-like environments the retained text needs and adds the article's own macro definitions that the retained text needs (\texttt{\textbackslash Z}, \texttt{\textbackslash tl}, \texttt{\textbackslash wh}), with extra wrapper packages (bm, bbm, dsfont, xspace, cancel, mathtools). The delivered numbering is the unit's own. Assets: no retained span contains a figure environment, a graphics-inclusion command, a PDF-page inclusion, a table or a TikZ picture; no image file is copied into this package, the package declares no asset dependency and no image file is redistributed with it. Licence: version arXiv:2605.18278v1 of this article is distributed under the Creative Commons Attribution 4.0 International licence, as recorded in the arXiv licence metadata for this exact version and shown on the abstract page https://arxiv.org/abs/2605.18278v1. A full rights notice is delivered at licenses/arxiv-2605.18278-CC-BY-4.0.txt. Characters: every retained excerpt is pure ASCII, so the delivered engine, XeLaTeX with its default Latin Modern text font, reports no missing character.
\begin{prop}
Both one-sided and two-sided Bratteli diagrams of infinite odometers $B_{IO}$ and $\tl B_{IO}$ are reducible, their corresponding tail equivalence relations $\mathcal{R}_{IO}$ and $\widetilde{\mathcal{R}}_{IO}$ are topologically transitive but not minimal. 
\end{prop}

\begin{proof}
For both diagrams $B_{IO}$ and $\tl B_{IO}$, from a vertex $i$ on level $n$ we will never reach vertex $j > i$ on any level below. Hence, both diagrams $B_{IO}$ and $\tl B_{IO}$ are reducible.

For the one-sided diagram $B_{IO}$, 
the set $\wh{X}_{B_{RS}}$ is dense in $X_{B_{IO}}$. The orbit a path $x \in X_{B_{IO}}$ is dense in $X_{B_{IO}}$ if and only if it $x$ is tail equivalent to a path from $X_{B_{RS}}$. Thus, the tail equivalence relation $\mathcal{R}_{IO}$ is topologically transitive but not minimal.

%any vertical path which passes through the vertex $0$ has a dense orbit, but the orbit of a vertical path which passes through the vertex $1$ will never visit the cylinder $[0]$ which corresponds to the vertex $0 \in V_0$.

For the two-sided diagram of infinite odometers $\tl B_{IO}$, there is no $i \in \Z$ such that $\wh{X}_{B_i}$ is dense in $X_{\tl B_{IO}}$. The tail equivalence class of any eventually vertical path is not dense, but we can find other paths that have dense orbits under the tail equivalence relation on the path space of $\tl B_{IO}$. 
First, note that any path $x \in X_{\tl B_{IO}}$ divides the set of all vertices of $\tl B_{IO}$ into three disjoint nonempty subsets: the set $V(x) = \{s(x_n)\}_{n = 0}^{\infty}$ of vertices through which $x$ passes, the set 
$$V_{l}(x) = \bigcup_{n \geq 0}\{w \in V_n : w < s(x_n)\}$$ 
%of vertices that are above or to the left of $x$, 
%such that $w \in V_n$ belongs to $V_l(x)$ if $w < s(x_n)$,
and the set 
$$V_{r}(x) = \bigcup_{n \geq 0}\{w \in V_n : w > s(x_n)\}.$$ 
%$V_r(x)$ 
%of the vertices %which are below or to the right of $x$.

Clearly, a path with a dense orbit cannot be vertical. Indeed, the paths from the tail equivalence class of a vertical path passing through the vertex $i$ on each level can reach only vertices $j \geq i$ on the levels above, i.e. the vertices from $V_r(x)$. Thus, they will never belong to a cylinder set which ends in a vertex from $V_l(x)$. In particular, they will never visit a cylinder set which corresponds to the vertex $i - 1$ on level $0$. 
     The slanted paths which start at some vertex $i \in V_0 = \Z$ and go all the time through the slanted edges also don't have dense orbits. Indeed, orbit of such a path will never reach a cylinder set which corresponds to the vertex $i+1$ on level $0$ or, in general, any other vertex from $V_r(x)$.
     
     The path $x$ which starts at a vertex $i \in V_0$ and passes in turn through vertical and slanted edges has a dense orbit. In other words, we consider the path which passes through the vertices $i, i, i - 1, i - 1, \ldots, i - k, i - k, \ldots$. 
     For simplicity, set $i = 0$, then we have $s(x_{2n}) = s(x_{2n + 1}) = -n$. Thus, $x$ passes through the vertex number $-n$ on levels $2n+1$ and $2n+2$. 
     %Note that the tower which corresponds to the vertex $i$ on level $l$ reaches the vertices $i, \ldots, i + (l - k)$ on the level $k < l$. Thus, the tower which corresponds to vertex $-n$ on the level $l = 2n$, reaches the vertices $-n, \ldots \tcr{compute}$
      Let $[(y_0, \ldots, y_n)]$ be any cylinder set that ends at the vertex $j^{(n+1)}$ on the level $n + 1$. If $j^{(n+1)} \in V_{l}(x)$, then we can use a vertical path that starts at $j^{(n+1)}$ to reach a vertex from $V(x)$. Thus, a point from the orbit of $x$ will belong to the cylinder $[(y_0, \ldots, y_m)]$. If $j^{(m+1)} \in V_{r}(x)$, we use the slanted path $sl(j^{(m+1)})$ that starts at $j^{(m+1)}$ to reach a vertex from $V(x)$. Such a path passes through the vertex $j^{(m+1)} - k$ on the level $m + 1 + k$.  If $j^{(m+1)} > 0$ then the path $sl(j^{(m+1)})$ passes through the vertex $0$ on the level $m + 1 + j^{(m+1)}$ and it passes though vertex $-n$ on the level $m + 1 + j^{(m+1)} + n$. Thus, for $n = j^{(m+1)} + m$ and $n = j^{(m+1)} + m - 1$, the slanted path $sl(j^{(m+1)})$ passes through the vertex which belongs to $V(x)$. If $j^{(m+1)} \leq 0$, we shift all the vertices uniformly by the same number $-r(x_m)$ to the right, so that now we have $r(x_m) = 0$, and consider only the part of $\tl B_{IO}$ which starts at the level $m+1$, Thus, we repeat the same proof as in the previous case, i.e. the proof for $j^{(m+1)} > 0$.  

      Note that the same proof works for any path $x \in X_{\tl B_{IO}}$ has infinitely many vertical and infinitely many slanted edges. In other words, $x \in X_{\tl B_{IO}}$ has a dense orbit if and only if $x$ is not tail equivalent to a vertical or to a slanted path. Such a point $x$ will ``intersect'' any slanted path that starts at a vertex from $V_r(x)$ and any vertical path that starts at a vertex from $V_{l}(x)$. Since $\tl B_{IO}$ contains only edges which are vertical or slanted that join a vertex $i$ to a vertex $i - 1$ on a level below, any path which passes through vertices from both $V_r(x)$ and $V_l(x)$ should pass through a vertex from $V(x)$. 
      %one cannot move from V_r on level n to V_l on level n+1.  Or from V_l above to V_r below.
      %some geometry here
      %for more explanation consider cases: V_l(x), V(x), V_r(x) on a level above and below, and if x has a vertical or a slanted edge 
\end{proof}

\end{document}
